% Generated by scripts/build_manuscript_latex.py from the Markdown master; do not edit.
\abstract{Generative foundation models that combine natural language with omics data must decide how biology reaches the language model: by translating it into the model's own tokens or by encoding it with a biological model and injecting the result. In this scoping review (searches to 9 August 2026) we included 55 records describing 51 independent works, 109 models and 585 input routes, decomposed the routes into a library of 22 data operations and charted rationales, supervision, dependence tests and cross-representation comparisons. Seventeen works are token-native, 16 encoder-mediated and 8 hybrid, and 9 use language to condition a biology model. Token-native works keep a median of 100 gene names per cell, and only 2 of the 51 works write measured values as numbers. Works that inject encoder outputs give the language model 1–256 vectors per biological object, and a third of the route steps in the median encoder-mediated work are opaque learned modules. Twenty-five works train or evaluate on text written by another large language model, and only 22 test whether outputs depend on the biological input. In 31 cross-culture comparisons in 20 works, the authors' own representation won 28. Opponents were mostly released, unadapted forms of the other culture, and opponents trained by the proposing team lost all 11 comparisons. The two cultures keep complementary parts of biological data, and the published comparisons cannot yet tell which serves a given task better. We propose a neutral comparison protocol and a reporting card for input representations.}

\keywords{large language models, foundation models, single-cell transcriptomics, multimodal learning, input representation, scoping review}
